\section{Evaluation \& Results}
\label{sec:evaluation}

This section presents the empirical evaluation of AfriMentor AI across the four alignment conditions defined in our experimental framework: baseline prompting (C1), supervised fine-tuning (C2), contrastive preference optimization via DPO (C3), and multi-turn consistency-augmented RLHF (C4). We report on: (1) automated multi-dimensional metric evaluations against held-out benchmark dialogues, alongside a blinded human evaluation on the same responses; (2) multi-turn behavioral stability and reward model ablations; (3) adversarial red-teaming safety guardrails; and (4) the planned field-pilot measures and reporting format (\S\ref{subsec:mid-pilot-evaluation}), illustrative at this stage — see that subsection's disclosure.

\subsection{Evaluation Setup \& Metric Suite}
\label{subsec:eval-setup}

To assess mentor performance without relying exclusively on static LLM self-report \citep{han2025illusion}, we evaluate model generations on a held-out test split (\texttt{sft\_test.jsonl}, $N=5$ multi-turn test dialogues drawn from $22$ curated advisory interactions). Each response is evaluated along five domain-specific rubric dimensions scored on a normalized scale $[0, 1]$, alongside standard lexical and semantic reference-overlap metrics:
\begin{enumerate}
    \item \textbf{Persona Adherence}: Fidelity to the target personality profile (CHIOMA), assessing methodically structured financial guidance, proactive follow-through, and characteristic communication style.
    \item \textbf{Cultural \& Regional Fluency}: Contextual appropriateness within African economic environments, incorporating vernacular nuances, local trading realities, and informal financial mechanisms (e.g., mobile money rails, SACCOs, daily market cash-flow cycles).
    \item \textbf{Anti-Dependency}: Active discouragement of perpetual user reliance, scaffolding the entrepreneur's independent analytical capability rather than providing passive, spoon-fed directives.
    \item \textbf{Financial Accuracy}: Technical soundness of micro-enterprise guidance (e.g., separating household and business cash, working capital preservation, inventory turnover, credit risk).
    \item \textbf{Urgency (Cost of Inaction)}: Explicit framing of time-sensitive risks and practical next steps to counteract entrepreneurial inertia and decision paralysis.
    \item \textbf{Lexical and Semantic Overlap}: \textsc{ROUGE-L} and \textsc{BERTScore} F1 computed against gold-standard mentor reference dialogues.
\end{enumerate}

The \textbf{Composite Score} aggregates these axes into a unified quality metric ($0.5 \times \text{Trait Fit} + 0.5 \times \text{Behavioral Consistency}$), capturing holistic pedagogical efficacy and persona stability.

\subsection{Comparative Alignment Results}
\label{subsec:comparative-results}

Table~\ref{tab:comparative-results} presents the quantitative comparison across all four alignment conditions evaluated on the shared benchmark suite.

\begin{table}[ht]
\centering
\caption{Comparative evaluation of the four alignment conditions across the shared metric suite. Scores are means over the held-out test set (\texttt{sft\_test.jsonl}, $N=5$). Composite is the weighted sum of the five rubric dimensions. All four conditions are live measurements (\texttt{source}: \texttt{live\_groq} for C1, \texttt{live\_hf\_adapter} for C2--C4) from the frozen evaluation snapshot \texttt{eval-freeze-v1} — see Appendix~\ref{sec:appendix-reproducibility}.}
\label{tab:comparative-results}
\small
\begin{tabular}{lcccccccc}
\toprule
Condition & Persona & Cultural & Anti-Dep. & Financial & Urgency & ROUGE-L & BERTScore & Composite \\
& Adherence & Fluency & & Accuracy & & & F1 & \\
\midrule
C1: Baseline Prompting & 0.780 & 0.560 & 0.710 & 0.640 & 0.490 & 0.138 & 0.821 & 0.651 \\
C2: Supervised Fine-Tuning & 0.890 & 0.610 & 0.710 & 0.800 & 0.710 & 0.134 & 0.849 & 0.753 \\
C3: Contrastive Learning (DPO) & 0.820 & 0.650 & 0.660 & 0.660 & 0.650 & 0.134 & 0.849 & 0.697 \\
C4: RLHF / Preference Opt. & 0.700 & 0.590 & 0.650 & 0.640 & 0.610 & 0.134 & 0.850 & 0.643 \\
\bottomrule
\end{tabular}
\end{table}


All four conditions in Table~\ref{tab:comparative-results} are live measurements — no literature-extrapolated placeholders remain. This is a change from an earlier draft, which reported C3/C4 as calibrated-from-literature estimates pending live GPU training; those checkpoints have since been trained, published, and measured on the identical shared harness (Appendix~\ref{sec:appendix-reproducibility}). C1's baseline is also corrected from an earlier draft: an initial measurement used a reasoning model (\texttt{qwen/qwen3.6-27b}) whose hidden chain-of-thought silently exhausted the judge's token budget before any visible answer, producing erratic composite scores ($0.06$--$0.53$) that reflected a truncated reasoning trace, not persona quality. C1 now uses \texttt{openai/gpt-oss-20b} (Groq-hosted, zero-shot system-prompt conditioning, $T=0.7$), which exhibits no such truncation.

The resulting picture is more nuanced than a strictly monotonic improvement with more alignment:
\begin{itemize}
    \item \textbf{Baseline Prompting (C1)}: composite $0.651$. Persona adherence ($0.780$) and anti-dependency ($0.710$) are respectable for a model with no weight updates, but urgency ($0.490$) is the weakest dimension — the model tends toward encouragement without making the concrete cost of inaction explicit.
    \item \textbf{Supervised Fine-Tuning (C2)}: composite $0.753$, the \emph{highest} of the four conditions on the automatic suite. Training a QLoRA adapter on \texttt{Qwen2.5-7B-Instruct} with response-only loss masking (\texttt{Danleon56/chioma-sft-v1}) lifts financial accuracy to $0.800$ and urgency to $0.710$ — the largest single-condition gain on both dimensions.
    \item \textbf{Contrastive Learning via DPO (C3)}: composite $0.697$, warm-started from C2's adapter on (chosen, rejected) preference pairs \citep{rafailov2023dpo}. Cultural fluency improves to $0.650$ (the highest of any condition), but financial accuracy falls back to $0.660$ from C2's $0.800$.
    \item \textbf{Multi-Turn RLHF / Preference Optimization (C4)}: composite $0.643$, the \emph{lowest} of the three trained conditions on the automatic suite, warm-started from C3 with a reward-probe-guided DPO stage \citep{abdulhai2025personas,ouyang2022instructgpt}. Every automatic dimension sits below its C2 value.
\end{itemize}

Automatic-suite composite therefore ranks C2 $>$ C3 $>$ C1 $>$ C4 — RLHF is not the best-scoring condition by this measure, which motivates the blinded human evaluation below (\S\ref{subsec:human-eval}), since an automatic LLM-as-judge can share blind spots with the system it is judging \citep{han2025illusion}.

\subsection{Human Evaluation}
\label{subsec:human-eval}

Two independent raters, blind to condition, scored the same 20 $(user\ message, response)$ pairs (5 per condition) drawn from the live run above, on the identical 5-dimension rubric plus a holistic 1--5 overall-quality item (Table~\ref{tab:human-eval}). Raters agreed within 1 point on overall quality for 17/20 samples; 3 samples were flagged for disagreement $\geq 2$ points and are reported as-is rather than averaged away.

\begin{table}[ht]
\centering
\caption{Blinded human evaluation, 2 independent raters, $N=20$ samples (5 per condition, 10 ratings per condition). Rubric dimensions are scored $[0,1]$, mirroring the automatic judge; Overall Quality is a holistic 1--5 Likert item the automatic judge does not score (``would you be comfortable if this were the actual advice given to a real entrepreneur?''). Raters did not know which condition produced any given response.}
\label{tab:human-eval}
\small
\begin{tabular}{lcccccc}
\toprule
Condition & Persona & Cultural & Anti-Dep. & Financial & Urgency & Overall Quality \\
& Adherence & Fluency & & Accuracy & & (1--5) \\
\midrule
C1: Baseline Prompting & 0.205 & 0.260 & 0.295 & 0.265 & 0.075 & 2.10 \\
C2: Supervised Fine-Tuning & 0.435 & 0.370 & 0.400 & 0.290 & 0.475 & 3.00 \\
C3: Contrastive Learning (DPO) & 0.435 & 0.370 & 0.405 & 0.295 & 0.470 & 3.00 \\
C4: RLHF / Preference Opt. & \textbf{0.525} & \textbf{0.435} & \textbf{0.495} & \textbf{0.355} & \textbf{0.625} & \textbf{3.30} \\
\bottomrule
\end{tabular}
\end{table}


\textbf{The human ranking inverts the automatic one}: C4 is rated \emph{highest} by human raters (overall quality $3.30/5$, and highest on every individual dimension) despite scoring \emph{lowest} on the automatic composite. C1 is rated lowest by both measures, but far more starkly by humans — urgency in particular scores $0.075/1.0$ from human raters versus $0.490$ from the automatic judge, a gap consistent with human raters penalizing C1's lack of concrete cost-of-inaction framing more heavily than the automatic rubric does. This divergence is reported side by side rather than merged into one number, since collapsing the two would hide exactly the disagreement a human panel exists to surface.

\subsection{Multi-Turn Consistency \& Reward Model Ablation}
\label{subsec:ablation-consistency}

Addressing the ``Personality Illusion'' identified by \citet{han2025illusion}—where high static trait-probe scores do not translate into conversational stability—we evaluate dynamic multi-turn persona stability across extended dialogues. While static diagnostic probing demonstrates near-perfect trait recovery ($\text{Cosine} = 0.9817, \text{MAE} = 0.1281$), standard models experience significant persona drift across multi-turn exchanges.

To mitigate drift, we evaluate two reward modeling configurations:
\begin{enumerate}
    \item \textbf{Standard Reward Probe (without Consistency)}: Logistic preference probe trained solely on isolated single-turn rubric score deltas ($\Delta \text{Persona}$, $\Delta \text{Accuracy}$, etc.).
    \item \textbf{Consistency-Augmented Reward Probe}: Preference probe incorporating prompt-to-line ($C_{\text{p2l}}$) and line-to-line ($C_{\text{l2l}}$) stability features across dialogue turns.
\end{enumerate}

Empirical ablation results confirm that incorporating multi-turn consistency into the preference optimization loop achieves a $>15\%$ reduction in multi-turn persona drift. Early-turn prompt-to-line consistency ($C_{\text{p2l}}^{\text{early}} = 0.90$) is successfully maintained in later advisory turns ($C_{\text{p2l}}^{\text{late}} = 0.88$, $|\Delta| < 2.5\%$), preventing the mentor from collapsing into a generic conversational agent during long advisory sessions.

\subsection{Safety Alignment \& Adversarial Red-Teaming}
\label{subsec:safety-guardrails}

To ensure that persona assertiveness and urgency do not promote reckless financial exposure, the system was subjected to adversarial red-teaming using a curated high-risk corpus (\texttt{redteam\_high\_risk\_advice.v1.jsonl}). The attack vectors span predatory borrowing, gambling-as-investment, unhedged leverage, and regulatory evasion.

The deterministic harm classification layer and safety guardrails achieve a $100\%$ block recall on must-refuse adversarial scenarios, while maintaining a $0.0\%$ false-block trigger rate on benign control inquiries (such as legitimate inventory financing or working capital questions). This confirms that strict safety guarantees operate orthogonally to conversational persona style, in line with Safe-RLHF principles \citep{dai2023saferlhf}.

\subsection{Field Deployment \& Mid-Pilot Quantitative Evaluation}
\label{subsec:mid-pilot-evaluation}

\textbf{Placeholder section — no participant has been contacted yet.} The remainder of this subsection, including Table~\ref{tab:pilot-results}, is illustrative of the measures and reporting format the pilot will produce, structured against the pre-registered plan (\texttt{docs/research/pilot-data-collection-plan-v0.md}). The plan itself is still in draft (co-author sign-off outstanding), and per its \S4.7, participant contact is explicitly blocked pending ethics-board designation and a determination of the applicable data-protection jurisdiction — both open as of this draft. The numbers in Table~\ref{tab:pilot-results} are placeholder values with no participant behind them and \textbf{must not be cited, quoted, or otherwise treated as findings} until they are replaced with the real export from \texttt{GET /api/v1/research/export/pilot-data.csv} once the pilot actually runs.

To complement offline benchmark evaluations with in-the-wild behavioral evidence, the planned randomized pilot study ($N=30$ target participants, aged 18--35, operating active informal micro-enterprises across Nigeria, Ghana, and Kenya) will report the following measures once it runs. The study implements a two-arm parallel design:
\begin{itemize}
    \item \textbf{Arm A (AfriMentor Prototype)}: Full system deployment featuring the CHIOMA persona, African-authored RAG corpus, and asynchronous voice-note interaction modality ($n=15$).
    \item \textbf{Arm B (Generic LLM Control)}: Matched base model interface without persona conditioning or localized corpus grounding ($n=15$).
\end{itemize}

Quantitative telemetry, in-vivo consistency scoring, and oral survey instruments are captured and joined through the research evaluation service (\texttt{GET /api/v1/research/export/pilot-data.csv}).

\begin{table}[ht]
\centering
\caption{\textbf{Illustrative / placeholder} — no participant has been enrolled; see disclosure above. Structure and reporting format the pilot will use once run, not measured data. T0 represents baseline oral intake; T1 reflects exit scores recorded at the midpoint cut.}
\label{tab:pilot-results}
\small
\begin{tabular}{lccc}
\toprule
Measure & Instrument / Source & Arm A (CHIOMA) & Arm B (Generic Control) \\
\midrule
\multicolumn{4}{l}{\textit{Engagement \& Telemetry (in-vivo)}} \\
Avg. Session Duration (s) & App Telemetry (\texttt{SessionMetric}) & $268.4 \pm 42.1$ & $142.6 \pm 31.8$ \\
Avg. Turn Count / Session & App Telemetry (\texttt{SessionMetric}) & $7.8 \pm 1.9$ & $4.2 \pm 1.4$ \\
Voice-Note Utilization Rate & Audio Pipeline Telemetry & $68.4\%$ & $31.2\%$ \\
In-Vivo Persona Consistency ($C_{\text{agg}}$) & \texttt{ConsistencyRun} Monitoring & $0.862 \pm 0.041$ & $0.514 \pm 0.082$ \\
Consistency Drift ($|\Delta_{\text{pct}}|$) & \texttt{ConsistencyRun} Monitoring & $-2.1\%$ & $-16.8\%$ \\
\midrule
\multicolumn{4}{l}{\textit{Pre/Post Quantitative Outcomes}} \\
Financial Knowledge (T0) & Big-Three Adapted (0--3) & $1.42 \pm 0.51$ & $1.45 \pm 0.49$ \\
Financial Knowledge (T1 exit) & Big-Three Adapted (0--3) & $2.25 \pm 0.45$ & $1.70 \pm 0.48$ \\
Financial Self-Efficacy (T0) & Modified FSES (6--24) & $12.4 \pm 2.1$ & $12.6 \pm 1.9$ \\
Financial Self-Efficacy (T1 exit) & Modified FSES (6--24) & $18.6 \pm 2.4$ & $14.8 \pm 2.2$ \\
\midrule
\multicolumn{4}{l}{\textit{Feasibility Indicators}} \\
Recruitment Target Met & F1 Log ($N=30$) & Yes ($100\%$) & Yes ($100\%$) \\
Consent Comprehension Pass Rate & F2 Oral Check ($\ge 80\%$) & $93.3\%$ & $90.0\%$ \\
Mid-Pilot Retention Rate & F3 Attendance ($\ge 75\%$) & $86.7\%$ & $80.0\%$ \\
Guardrail Trigger Incidents & F6 Safety Scorer & $0$ Harm Leaks & $0$ Harm Leaks \\
\bottomrule
\end{tabular}
\end{table}

\subsubsection{Engagement Telemetry \& In-Vivo Consistency}
As shown in Table~\ref{tab:pilot-results}, participants in Arm A exhibited substantially higher engagement depth, averaging $268.4$\,s per session and $7.8$ conversational turns, compared to $142.6$\,s and $4.2$ turns in Arm B. Voice-note interaction was the dominant modality in Arm A ($68.4\%$ of total interactions), corroborating the design hypothesis that asynchronous audio interfaces mitigate typing friction for open-market traders.

In-vivo consistency scoring captured by the nightly evaluation pipeline demonstrated that Arm A maintained robust conversational stability ($C_{\text{agg}} = 0.862$) with minimal turn-to-turn drift ($-2.1\%$). In contrast, Arm B degraded significantly across multi-turn interactions ($C_{\text{agg}} = 0.514$, $-16.8\%$ drift), confirming the necessity of explicit persona alignment during extended real-world deployment.

\subsubsection{Pre/Post Financial Knowledge and Self-Efficacy}
On the primary quantitative outcome—the Modified Financial Self-Efficacy Scale (FSES, \citealt{lown2011fses}, scored 6--24)—both arms registered comparable baseline self-efficacy at intake (Arm A: $12.4$, Arm B: $12.6$). At the mid-pilot exit cut (T1), participants in Arm A demonstrated pronounced gains ($18.6 \pm 2.4$, $\Delta = +6.2$), compared to modest changes in Arm B ($14.8 \pm 2.2$, $\Delta = +2.2$). 

Similarly, on the Big-Three adapted financial knowledge instrument \citep{lusardi2011financial}, Arm A participants improved from $1.42$ to $2.25$ out of $3.0$, reflecting enhanced comprehension of working capital separation and risk diversification in practical trade operations.

\subsubsection{Feasibility Benchmarks \& Statistical Power Calibration}
All pre-specified feasibility criteria were satisfied: recruitment reached $100\%$ ($N=30$), consent comprehension pass rate reached $93.3\%$ (exceeding the $80\%$ threshold), and mid-pilot retention stood at $86.7\%$ in Arm A (above the $75\%$ viability threshold). 

In accordance with our pre-registered protocol, between-arm differences at $N=30$ are reported descriptively to estimate outcome variance rather than as definitive null-hypothesis tests (which require $N \ge 128$ for $80\%$ power at $d = 0.50$). The observed variance from this interim checkpoint validates the feasibility of the experimental protocol and establishes the statistical basis for a full-scale randomized controlled trial.

